\documentclass[../../../include/open-logic-section]{subfiles}

\begin{document}

\olfileid[id]{sth}{ordinals}{replacement}
\olsection{Penggantian}

Dalam \olref[sth][ordinals][vn]{sec}, kita memotivasi pengenalan ordinal
dengan menyarankan bahwa kita dapat memperlakukannya sebagai tipe urutan,
yakni sebagai wakil kanonik bagi pengurutan baik. Agar hal itu berhasil, kita
perlu membuktikan bahwa \emph{setiap pengurutan baik isomorfik dengan suatu
ordinal}. Ini akan memungkinkan kita mendefinisikan $\ordtype{A, <}$ sebagai
ordinal $\alpha$ sedemikian sehingga $\tuple{A, <} \isomorphic \alpha$.

Sayangnya, kita \emph{tidak dapat} membuktikan hasil yang diinginkan hanya
dengan aksioma-aksioma yang telah diperkenalkan sejauh ini. (Kita akan
melihat alasannya dalam \olref[replacement][strength]{sec}, tetapi untuk saat
ini pokoknya ialah: kita tidak dapat melakukannya.) Kita memerlukan suatu
gagasan baru, yakni:

\begin{axiom}[Skema Penggantian]
Untuk sembarang formula $\phi(x, y)$, berikut ini merupakan suatu aksioma:
\begin{quote}
	untuk sembarang $A$, jika $(\forall x \in A)\lexists![y][\phi(x,y)]$,
	maka $\Setabs{y}{(\exists x \in A)\phi(x,y)}$ ada.
\end{quote}
\end{axiom}
\noindent
Seperti Pemisahan, ini merupakan suatu skema: skema tersebut menghasilkan
tak berhingga banyaknya aksioma, satu bagi setiap formula $\phi$ yang berbeda.
Skema ini juga dapat (dan biasanya memang) dituliskan sebagai berikut:

\begin{defish}
Untuk sembarang formula $\phi(x,y)$ yang tidak memuat ``$B$'', berikut ini
merupakan suatu aksioma:
\[
\forall A[(\forall x \in A)\lexists![y][\phi(x,y)] \lif \exists B\forall y (y \in B \liff (\exists x \in A)\phi(x,y))]
\]
\end{defish}

Namun, ketika pertama kali dijumpai, rumus ini cukup berbelit-belit.
Konsekuensi singkat Penggantian berikut mungkin mengungkapkan gagasan
intuitif yang sedang kita gunakan dengan \emph{lebih jelas}:\footnote{Suatu
suku adalah ungkapan yang menunjuk tepat satu objek (pada setiap pengisian
variabel bebasnya). Sebagai contoh, ``$\emptyset$'' merupakan suku yang
menunjuk himpunan kosong; ``$\{x\}$'' merupakan suku yang menunjuk singleton
$x$ (apa pun $x$ itu); ``$x \cup y$'' merupakan suku yang menunjuk gabungan
$x$ dan $y$ (apa pun keduanya).}

\begin{cor}
Untuk sembarang suku $\tau(x)$ dan sembarang himpunan $A$, himpunan berikut
ada:
\[
	\Setabs{\tau(x)}{x \in A} = \Setabs{y}{(\exists x \in A)y = \tau(x)}.
\]
\end{cor}

\begin{proof}
Karena $\tau$ merupakan suatu \emph{suku}, $\forall x
\lexists![y][\tau(x) = y]$. Terlebih lagi, $(\forall x \in A)
\lexists![y][\tau(x) = y]$. Jadi, $\Setabs{y}{(\exists x \in A)\tau(x) =
y}$ ada menurut Penggantian.
\end{proof}
\noindent
Hal ini menunjukkan bahwa ``Penggantian'' merupakan nama yang tepat bagi
aksioma tersebut: dengan diberikannya suatu himpunan $A$, Anda dapat
membentuk himpunan baru, $\Setabs{\tau(x)}{x \in A}$, dengan mengganti setiap
anggota $A$ dengan citranya di bawah~$\tau$. Memang, mengikuti notasi bagi
citra suatu himpunan di bawah suatu fungsi, kita dapat menulis
$\funimage{\tau}{A}$ untuk $\Setabs{\tau(x)}{x \in A}$.

Namun, yang penting, $\tau$ merupakan suatu \emph{suku}. Suku itu tidak harus
merupakan (nama bagi) suatu \emph{fungsi}, dalam pengertian
\olref[sfr][fun][rel]{sec}, yakni suatu himpunan tertentu dari pasangan
terurut. Lagi pula, jika $f$ merupakan suatu fungsi (dalam pengertian itu),
maka himpunan $\funimage{f}{A} = \Setabs{f(x)}{x \in A}$ hanyalah suatu
himpunan bagian tertentu dari $\ran{f}$, dan keberadaannya sudah dijamin
hanya dengan menggunakan aksioma-aksioma~$\Zminus$.\footnote{Cukup
pertimbangkan $\Setabs{y \in \bigcup \bigcup f}{(\exists x \in A)y =
f(x)}$.} Sebaliknya, Penggantian merupakan tambahan yang \emph{kuat} bagi
aksioma-aksioma kita, sebagaimana akan kita lihat dalam
\olref[sth][replacement][]{chap}.


\end{document}
